%%APPEND title="Ordinary covers yield homotopy pushouts, with one direction of excision"
The question left on ms.~76, whether $W$-asphericity of the intersection and of the union forces that of the two pieces without the converse half of W(7 bis), has an affirmative answer for every fundamental localizer satisfying the relative axiom; such a localizer is strongly saturated. In outline: let $K=U\cap V$ and $X=U\cup V$ be $W$-aspherical. If $K$ is empty then $\emptyset\to e$ is in $W$ and $W$ consists of all functors. Otherwise choose an object $k$ of $K$; the functor $e\to K$ it defines is in $W$. The forward implication, applied to the map from the span $U\leftarrow e\to V$ to the span $U\leftarrow K\to V$, identifies $X$ in the localized category with the wedge of $U$ and $V$ at $k$, which is therefore $W$-aspherical. The arrow $U\to e$ is a retract of the arrow from the wedge to $e$, and strong saturation gives $U\to e$ in $W$; likewise for $V$. The retract step is the one written out in the commentary at ms.~272 (``The retract argument is valid at the level of objects''). Only the forward implication is used.

%%APPEND title="A finite counterexample to the asserted reverse excision"
Two complements. A smaller category serves the same purpose: the face category of a triangulation of the presentation complex of $\langle a,b\mid a^5b^{-3},\,b^3(ab)^{-2}\rangle$ with one cone point in each relation cell has 35 objects and 159 arrows, is connected and integrally acyclic, and maps onto $A_5$ on fundamental groups. And the failure does not depend on a badly placed sieve: it persists for a cover in which the inclusion of the intersection is a Dwyer map, obtained from the mapping cylinder of $A\to e$ in place of a cone; that verification is not reproduced here.

%%APPEND title="What survives of the cube principle"
The exact cost of the reverse axiom is now known. A localizer satisfying the reverse implication of W(7 bis) contains every functor inducing isomorphisms on integral homology, and the class of those functors is the least such localizer; a localizer satisfying it and different from the class of all functors makes no sphere aspherical and never has enough intrinsic fibrations. The proofs are in the commentary on Chapter~XVI, ms.~27. The conditional propositions of ms.~81--83, which assume W(7 bis) or W(7), are therefore statements about homological localizers and not about $W_\infty$, also when the inclusions involved are Dwyer maps.
